\documentclass[10pt,a4paper]{amsart}
\usepackage[margin=19mm]{geometry}
\usepackage{amsmath,amssymb,mathtools}
\usepackage[scr=rsfs,cal=euler]{mathalfa}
\usepackage{xfrac}
\usepackage[unicode,hidelinks,bookmarks=false]{hyperref}
\newcommand{\ie}{i.e.\ }
\newcommand{\cf}{cf.\ }
\newcommand{\resp}{respectively\ }
\newcommand{\Subsection}{\subsection}
\providecommand{\nobreakdash}{\nobreak\hbox{-}}
\setlength{\emergencystretch}{2em}
\multlinegap0pt
\usepackage{bm}
\usepackage{bbm}
\usepackage{dsfont}
\usepackage{xspace}
\usepackage{cancel}
\usepackage{mathtools}
\usepackage{amsthm}
\makeatletter
\@ifundefined{theorem}{\newtheorem{theorem}{Theorem}}{}
\@ifundefined{proposition}{\newtheorem{proposition}[theorem]{Proposition}}{}
\makeatother
\newcommand{\KG}{\mathcal{K}_{G}}
\newcommand{\On}{\mathrm{O}(n)}
\title[The $L\_{p}$ Dual Minkowski Problem for Group-Invariant Conve]{The $L\_{p}$ Dual Minkowski Problem for Group-Invariant Convex Bodies}
\author{Junjie Shan}
\date{Source: arXiv:2507.13233v2 (CC BY 4.0)}
\begin{document}
\maketitle

\noindent\emph{Source and licence.} The $L_{p}$ Dual Minkowski Problem for Group-Invariant Convex Bodies. Version arXiv:2507.13233v2 (CC BY 4.0), distributed under the Creative Commons Attribution 4.0 International licence, as recorded in the arXiv licence metadata for this exact version and shown on the abstract page https://arxiv.org/abs/2507.13233v2. Full rights notice: licenses/arxiv-2507.13233-CC-BY-4.0.txt.\par\smallskip

\noindent\textit{Editorial scope.} Counted unit: the Proposition whose source label is \texttt{G=osym} in the article (source0.tex lines 664--671, source label \texttt{G=osym}), delivered together with the article's own complete original proof of it (source0.tex lines 673--730, one \texttt{proof} environment of 58 lines). No mathematics is written, repaired or re-derived: statement and proof are the article's own text line for line. Required context retained uncounted: none. Every definition, display and label of those lines is retained unchanged. Omitted from this package and not redistributed: 1--663, 672--672, 731--2385. Nothing inside a retained excerpt is omitted or altered: every retained body is delivered exactly as the article's lines are written, so no omission receipt of any kind accompanies this package. Citations: the retained excerpts carry no citation command at all, so no bibliography entry of the article is transcribed and no reference is redistributed. Reference adaptation: none was needed and none was made; every reference inside the delivered unit resolves inside it, so no unresolved reference or citation is produced. Printed slips kept literal: no printed word, symbol, hypothesis or number of the counted statement or of its proof was altered. Layout and class: the article's own class is \texttt{amsart} with packages including mathrsfs, amsfonts, amsmath, amsthm, amssymb, natbib, graphicx, comment, diagbox, tikz, booktabs, makecell, float, multirow; the wrapper is the lane's pinned \texttt{amsart} editorial wrapper at 10pt on A4, which restates the theorem-like environments the retained text needs and adds the article's own macro definitions that the retained text needs (\texttt{\textbackslash KG}, \texttt{\textbackslash On}), with extra wrapper packages (bm, bbm, dsfont, xspace, cancel, mathtools). The delivered numbering is the unit's own. Assets: no retained span contains a figure environment, a graphics-inclusion command, a PDF-page inclusion, a table or a TikZ picture; no image file is copied into this package, the package declares no asset dependency and no image file is redistributed with it. Licence: version arXiv:2507.13233v2 of this article is distributed under the Creative Commons Attribution 4.0 International licence, as recorded in the arXiv licence metadata for this exact version and shown on the abstract page https://arxiv.org/abs/2507.13233v2. A full rights notice is delivered at licenses/arxiv-2507.13233-CC-BY-4.0.txt. Characters: every retained excerpt is pure ASCII, so the delivered engine, XeLaTeX with its default Latin Modern text font, reports no missing character.
	\begin{proposition}\label{G=osym}
		Let $G$ be a subgroup of the orthogonal group $\On$. 
		The following conditions are equivalent:
		\begin{enumerate}
			\item $\KG=\mathcal{K}_{e}^{n}$.
			\item $G = \{\pm I\}$.
		\end{enumerate}
	\end{proposition}

	\begin{proof}
		We prove both directions of the equivalence.
		
		\noindent\textbf{(2 $\Rightarrow$ 1):} 
		Assume $G = \{\pm I\}$. 
		If $K$ is $G$-invariant, then for $g = -I \in G$, we have $(-I)K = -K = K$. 
			Thus $K$ is origin-symmetric.
		 If $K$ is origin-symmetric ($K = -K$), then
				$IK = K$, 
				$(-I)K = -K = K$,
			so $K$ is $G$-invariant.
		Hence the $G$-invariant convex bodies are precisely the origin-symmetric convex bodies.
		
		\noindent\textbf{(1 $\Rightarrow$ 2):} 
		Let $T: \mathbb{R}^n \to \mathbb{R}^n$ be a linear transformation such that for every $v \in \mathbb{R}^n$, either $Tv = v$ or $Tv = -v$. Then $T = \pm I$. Indeed, let $\{e_1,\ldots,e_n\}$ be an orthonormal basis. For each $i$, we have
		\begin{equation}\label{1}
			Te_i = \lambda_i e_i \quad \text{with } \lambda_i = \pm1.
		\end{equation}
	  For $i \neq j$, $T(e_{i}+e_{j})=\lambda_i e_i+\lambda_j e_j=\mu(e_{i}+e_{j})$
		for some $\mu = \pm 1$, implying $\lambda_i = \lambda_j$. Thus all $\lambda_i$ are equal. By \eqref{1} we have $T=\pm I$. 
	
		Assume condition (1) holds. We show that $G = \{\pm I\}$.
	If there exists $g \in G$ with $g \neq \pm I$,  we construct a convex body that is origin-symmetric but not $G$-invariant.
		Since $g \neq \pm I$, there exists a unit vector $v \in \mathbb{R}^n$ such that $gv \neq \pm v$. 
		Assume $v = e_1$ after a change of coordinates.
		
		Consider the ellipsoid (an origin-symmetric convex body)
		\begin{equation}\label{elliK}
		K = \{ x \in \mathbb{R}^n : x^T A x \leq 1 \},
		\end{equation}
		where $A = \operatorname{diag}(1, 2, \dots, 2)$ is a positive definite matrix. 
		Under the action of $g$, we have:
		\begin{align}\label{elligK}
		gK &= \{ gx \in \mathbb{R}^n : x^T A x \leq 1 \}\notag\\
		&=\{ y \in \mathbb{R}^n : (g^{-1}y)^T A (g^{-1}y) \leq 1 \}\notag\\
		&=\{ y \in \mathbb{R}^n : y^T (g A g^T) y \leq 1 \}.
		\end{align}
		
		Now compute the quadratic forms:
		\[
		e_1^T A e_1 = 1,
		\]
		\[
		e_1^T (g^{T} A g) e_1 = (g e_1)^T A (g e_1) = w^T A w, \quad \text{where} \quad w = g e_1.
		\]
		Since $g$ is orthogonal, $w$ is a unit vector. As $g e_1 \neq \pm e_1$, we have $w \neq \pm e_1$. 
	Let $w = (w_1, w_2, \dots, w_n)$ with $w_1^2 < 1$ and $\sum_{i=2}^n w_i^2 = 1 - w_1^2 > 0$. Then
		\[
		w^T A w =  w_1^2 + 2 \sum_{i=2}^n w_i^2 =  w_1^2 + 2 (1 - w_1^2)=2 - w_1^2 > 1  = e_1^T A e_1.
		\]
	 Thus $g^{T} A g \neq A$, since $g^{T}=g^{-1}$, so $g A g^{T} \neq A$. Then we have  $gK \neq K$ by \eqref{elliK} and \eqref{elligK}, contradicting 
	 $(1)$.
		Therefore, if $g\in G$, $g =  I$ or $g=-I$.
		
		 So we have $G \subset \{\pm I\}$. 
		If $G = \{I\}$, then all convex bodies (including non-origin-symmetric ones) are $G$-invariant, 
		contradicting  $(1)$. Thus $G = \{\pm I\}$.
	\end{proof}

\end{document}
